% Preámbulo modular para presentaciones Beamer (Modelación Estadística)
% Diseño y paleta institucional del Tecnológico de Monterrey

\documentclass[aspectratio=169,xcolor={dvipsnames,table}]{beamer}

\usepackage[utf8]{inputenc}
\usepackage[spanish,mexico]{babel}
\usepackage{amsmath,amssymb,amsfonts,mathtools}
\usepackage{graphicx}
\usepackage{listings}
\usepackage{booktabs}
\usepackage{tikz}
\usetikzlibrary{matrix,arrows,decorations.pathmorphing,shapes.geometric,positioning}

% Paleta Institucional Tecnológico de Monterrey
\definecolor{TecAzul}{HTML}{0039A6}       % Azul TEC: color primario institucional
\definecolor{TecAzulOscuro}{HTML}{1A2E51} % Azul marino oscuro
\definecolor{TecRojo}{HTML}{EC2661}       % Rojo de acento (paleta miTec)
\definecolor{TecGris}{HTML}{646464}       % Gris neutro para comentarios y subtítulos
\definecolor{TecFondoBloque}{HTML}{F4F6F9}% Fondo suave para bloques

% Configuración de Tema Beamer
\usetheme{Boadilla}
\usecolortheme[named=TecAzulOscuro]{structure}
\setbeamercolor{alerted text}{fg=TecRojo}
\setbeamercolor{palette primary}{bg=TecAzulOscuro,fg=white}
\setbeamercolor{palette secondary}{bg=TecAzul,fg=white}
\setbeamercolor{palette tertiary}{bg=TecAzulOscuro!80!black,fg=white}
\setbeamercolor{title}{fg=TecAzulOscuro,bg=TecFondoBloque}
\setbeamercolor{frametitle}{fg=TecAzulOscuro,bg=TecFondoBloque}

% Bloques personalizados elegantes
\setbeamercolor{block title}{bg=TecAzulOscuro,fg=white}
\setbeamercolor{block body}{bg=TecFondoBloque,fg=black}
\setbeamercolor{block title alerted}{bg=TecRojo,fg=white}
\setbeamercolor{block body alerted}{bg=TecRojo!8,fg=black}
\setbeamercolor{block title example}{bg=TecAzul,fg=white}
\setbeamercolor{block body example}{bg=TecAzul!8,fg=black}

% Redondear bordes y añadir sombra sutil
\setbeamertemplate{blocks}[rounded][shadow=true]
\setbeamertemplate{navigation symbols}{}

% Pie de página limpio con número de diapositiva
\setbeamertemplate{footline}{
  \leavevmode%
  \hbox{%
  \begin{beamercolorbox}[wd=.7\paperwidth,ht=2.25ex,dp=1ex,left]{author in head/foot}%
    \hspace*{2ex}\insertshortauthor~~(\insertshortinstitute) \hfill \insertshorttitle
  \end{beamercolorbox}%
  \begin{beamercolorbox}[wd=.3\paperwidth,ht=2.25ex,dp=1ex,right]{date in head/foot}%
    \insertshortdate{}\hspace*{2em}
    \insertframenumber{} / \inserttotalframenumber\hspace*{2ex}
  \end{beamercolorbox}}%
  \vskip0pt%
}

% Configuración de Listings de Python para visualización pedagógica de código
\lstset{
    language=Python,
    basicstyle=\ttfamily\footnotesize,
    keywordstyle=\color{TecAzulOscuro}\bfseries,
    stringstyle=\color{TecRojo},
    commentstyle=\color{TecGris}\itshape,
    showstringspaces=false,
    breaklines=true,
    frame=lines,
    rulecolor=\color{TecAzulOscuro!30},
    backgroundcolor=\color{TecFondoBloque},
    aboveskip=0.5em,
    belowskip=0.5em,
    literate={á}{{\'a}}1 {é}{{\'e}}1 {í}{{\'i}}1 {ó}{{\'o}}1 {ú}{{\'u}}1
             {Á}{{\'A}}1 {É}{{\'E}}1 {Í}{{\'I}}1 {Ó}{{\'O}}1 {Ú}{{\'U}}1
             {ñ}{{\~n}}1 {Ñ}{{\~N}}1 {¿}{{?`}}1 {¡}{{!`}}1
}

% Metadatos institucionales por defecto
\title[Modelación Estadística]{Modelación Estadística para Ciencia de Datos e Ingeniería}
\author[J. Castillo Colmenares]{Juliho David Castillo Colmenares}
\institute[Tec de Monterrey]{Tecnológico de Monterrey \\ Escuela de Ingeniería y Ciencias}
\date{\today}


\title{Medidas de Dispersión}
\subtitle{Sección 1.3 — Modelación Estadística para Ciencia de Datos}
\author[J. Castillo Colmenares]{Juliho Castillo Colmenares}
\institute{Tecnológico de Monterrey}
\date{}

\begin{document}

% Slide 1: Portada
\begin{frame}[plain]
  \titlepage
\end{frame}

% Slide 2: Hoja de Ruta e Indicador de Posición
\begin{frame}{Hoja de Ruta — Capítulo 01: Estadística Descriptiva}
  ¿Dónde estamos en nuestro recorrido por la estadística básica? \pause
  \begin{itemize}
    \item \textbf{01.01 Introducción a la estadística descriptiva} \emph{(Completado)} \pause
    \item \textbf{01.02 Medidas de tendencia central} \emph{(Completado)} \pause
    \item \color{TecRojo}\textbf{01.03 Medidas de dispersión \emph{(Hoy)}}\color{black}
  \end{itemize}
  \vspace{1em}
  \begin{block}{Objetivo de la Sección 1.3}
    Cuantificar qué tan dispersos, variables o concentrados están nuestros datos alrededor de su centro, comprendiendo la corrección de Bessel y las medidas robustas intercuartílicas.
  \end{block}
\end{frame}

% Slide 3: El Dilema de las Dos Muestras (Pregunta Gatillo)
\begin{frame}{El Dilema de las Dos Muestras}
  \begin{alertblock}{Situación de Latencia en Servidores (Pregunta Gatillo)}
    Un ingeniero de confiabilidad (SRE) monitorea los tiempos de respuesta de dos arquitecturas de servidores en la nube. Toma una muestra de 5 peticiones (en milisegundos) de cada una:
    \vspace{0.4em}
    \begin{itemize}
      \item \textbf{Arquitectura A:} $48, 49, 50, 51, 52 \text{ ms} \implies \bar{X}_A = 50.0 \text{ ms}$
      \item \textbf{Arquitectura B:} $10, 30, 50, 70, 90 \text{ ms} \implies \bar{X}_B = 50.0 \text{ ms}$
    \end{itemize}
    \vspace{0.4em}
    \begin{enumerate}
      \item Si ambos sistemas tienen exactamente la misma media ($\bar{X} = 50 \text{ ms}$), ¿son igual de confiables y estables en producción? \pause
      \item ¿Por qué la media aritmética y la mediana son totalmente ciegas ante la variabilidad del sistema?
    \end{enumerate}
  \end{alertblock}
\end{frame}

% Slide 4A: Conclusiones del Dilema (Análisis)
\begin{frame}{Conclusiones del Dilema — Análisis}
  \begin{itemize}
    \item \textbf{La Ceguera del Promedio:} La media únicamente indica dónde se equilibra el centro de masa ($\bar{X} = 50 \text{ ms}$), pero oculta completamente la oscilación de los extremos. \pause
    \item \textbf{Inestabilidad Operativa:} La Arquitectura A es consistente y predecible. La Arquitectura B es caótica; un usuario puede experimentar $10 \text{ ms}$ o un retraso severo de $90 \text{ ms}$. \pause
    \item \textbf{El Rol de la Dispersión:} Para caracterizar una distribución con rigor en Ciencia de Datos, el centro (\emph{Ubicación}) siempre debe ir acompañado de una métrica de dispersión (\emph{Escala}).
  \end{itemize}
\end{frame}

% Slide 4B: Conclusiones del Dilema (Respuestas Formales)
\begin{frame}{Conclusiones del Dilema — Respuestas}
  \begin{block}{Respuestas Formales al Dilema}
    \begin{enumerate}
      \item \textbf{¿Son igual de confiables?} Absolutamente no. La Arquitectura A ofrece una experiencia predecible, mientras que la B presenta una variabilidad que degrada la calidad del servicio. \pause
      \item \textbf{¿Por qué son ciegas la media y mediana?} Porque fueron diseñadas matemáticamente solo para ubicar la posición del centro, no para medir la distancia o separación de los datos entre sí.
    \end{enumerate}
  \end{block}
\end{frame}

% Slide 5: El Rango y la Desviación Media Absoluta (MAD)
\begin{frame}{Primeros Intentos: Rango y Desviación Media Absoluta}
  ¿Cómo podríamos medir la separación de los datos algebraicamente? \pause
  \begin{itemize}
    \item \textbf{El Rango ($R$):} Es la diferencia entre el dato máximo y el mínimo:
      \begin{align}
        R = X_{\max} - X_{\min}
      \end{align}
      \small \emph{Limitación:} Depende exclusivamente de dos números e ignora el 100\% del comportamiento intermedio. \pause
    \item \textbf{Desviación Media Absoluta (MAD):} Promedia el valor absoluto de las distancias al centro:
      \begin{align}
        \text{MAD} = \dfrac{1}{N} \sum_{i=1}^{N} |X_i - \bar{X}|
      \end{align}
      \small \emph{Limitación:} El valor absoluto $|x|$ genera esquinas no diferenciables en cálculo, dificultando la optimización matemática y el álgebra lineal.
  \end{itemize}
\end{frame}

% Slide 6: La Varianza Poblacional vs Muestral
\begin{frame}{La Varianza y la Corrección de Bessel}
  Para obtener una función diferenciable y tratable analíticamente, elevamos al cuadrado las desviaciones respecto a la media: \pause
  
  \begin{columns}[T]
    \begin{column}{0.48\textwidth}
      {\large \color{TecAzulOscuro}\textbf{Varianza Poblacional ($\sigma^2$)}}
      \vspace{0.2em}
      \begin{align}
        \sigma^2 = \dfrac{1}{N} \sum_{i=1}^{N} (X_i - \mu)^2
      \end{align}
      \small Promedio exacto del cuadrado de las desviaciones al conocer todo el universo $N$.
    \end{column}
    
    \begin{column}{0.48\textwidth}
      {\large \color{TecRojo}\textbf{Varianza Muestral ($S^2$)}}
      \vspace{0.2em}
      \begin{align}
        S^2 = \dfrac{1}{n - 1} \sum_{i=1}^{n} (X_i - \bar{X})^2
      \end{align}
      \small Se divide entre $n-1$ (\textbf{Corrección de Bessel}) para evitar sesgo al estimar $\sigma^2$.
    \end{column}
  \end{columns}
  \vspace{0.5em}
  \pause
  \begin{block}{¿Por qué $n-1$ grados de libertad?}
    Al usar $\bar{X}$ en lugar del verdadero $\mu$, los datos pierden 1 grado de libertad porque la suma $\sum(X_i - \bar{X})$ está forzada a ser 0. Dividir entre $n-1$ compensa que $\bar{X}$ está más cerca de la muestra que del parámetro real.
  \end{block}
\end{frame}

% Slide 7: Desviación Estándar y Coeficiente de Variación
\begin{frame}{Desviación Estándar ($S$) y Coeficiente de Variación ($CV$)}
  Dado que la varianza está en unidades al cuadrado ($\text{ms}^2$, $\text{pesos}^2$), tomamos la raíz cuadrada positiva para volver a las unidades originales de medición:
  \begin{align}
    S = \sqrt{S^2} = \sqrt{\dfrac{1}{n-1} \sum_{i=1}^{n} (X_i - \bar{X})^2}
  \end{align}
  \pause
  \begin{block}{El Coeficiente de Variación ($CV$)}
    ¿Cómo comparamos la dispersión del salario en dólares frente a la estatura en centímetros? El $CV$ normaliza la desviación respecto a la media en porcentaje:
    \begin{align}
      CV = \left( \dfrac{S}{\bar{X}} \right) \times 100\%
    \end{align}
    Un $CV > 30\%$ indica alta heterogeneidad y fuerte dispersión relativa.
  \end{block}
\end{frame}

% Slide 8: Medida Robusta — El Rango Intercuartílico (IQR)
\begin{frame}{Medidas Robustas: El Rango Intercuartílico ($IQR$)}
  Al igual que la media es vulnerable a valores extremos (\emph{outliers}), la desviación estándar $S$ eleva al cuadrado cualquier anomalía, magnificando su impacto. \pause
  \begin{itemize}
    \item \textbf{Los Cuartiles ($Q_1, Q_2, Q_3$):} Dividen la muestra ordenada en 4 partes del $25\%$ de los datos cada una ($Q_2$ es la Mediana). \pause
    \item \textbf{El Rango Intercuartílico ($IQR$):} Mide la dispersión del $50\%$ central de las observaciones:
  \end{itemize}
  \begin{align}
    IQR = Q_3 - Q_1
  \end{align}
  \pause
  \begin{alertblock}{Regla de Detección de Outliers (Tukey)}
    Se considera un valor atípico severo a cualquier observación por fuera de las cercas:
    \begin{align}
      \text{Límite Inferior} = Q_1 - 1.5 \cdot IQR, \quad \text{Límite Superior} = Q_3 + 1.5 \cdot IQR
    \end{align}
  \end{alertblock}
\end{frame}

% Slide 9A: Laboratorio en Python (Comparación de Varianzas y Bessel)
\begin{frame}{Laboratorio en Python: Varianza y Corrección de Bessel}
  Cálculo computacional demostrando la diferencia entre \texttt{ddof=0} (población) y \texttt{ddof=1} (muestra).
  \\ Script: \texttt{\footnotesize presentaciones/code/01\_estadistica\_descriptiva/01.03\_range\_variance\_stdev.py}
  
  \vspace{0.4em}
  {\lstset{basicstyle=\ttfamily\tiny, aboveskip=0.2em, belowskip=0.2em}
  \lstinputlisting[firstline=4,lastline=21]{../../code/01_estadistica_descriptiva/01.03_range_variance_stdev.py}
  }
\end{frame}

% Slide 14: Laboratorio en Python (Parte 2: Verificación con SciPy/Stats)
\begin{frame}[fragile]{Verificación en Python con SciPy (2/2)}
  \footnotesize
  Validación industrial de la Curtosis y Asimetría con la librería estándar de ciencia de datos (\texttt{01.03\_range\_variance\_stdev.py}):

  \vspace{0.2em}
  {\lstset{basicstyle=\ttfamily\tiny, aboveskip=0.2em, belowskip=0.2em}
  \lstinputlisting[firstline=24,lastline=39]{../../code/01_estadistica_descriptiva/01.03_range_variance_stdev.py}
  }
\end{frame}

% Slide 9C: Laboratorio en Python (Salida Verificada en Consola)
\begin{frame}{Laboratorio en Python: Verificación en Consola}
  Resultados empíricos obtenidos al ejecutar \texttt{01.03\_medidas\_dispersion.py} en la terminal:
  \vspace{0.6em}

  \begin{block}{\footnotesize Salida de Consola — Parte 1: Corrección de Bessel}
    {\fontsize{6.5pt}{7.5pt}\selectfont
    \texttt{--- Comparación de Muestras (Media vs Dispersión) ---} \\
    \texttt{Muestra A: Media = 50.0 ms | Desv. Estándar (S) = 1.58 ms} \\
    \texttt{Muestra B: Media = 50.0 ms | Desv. Estándar (S) = 31.62 ms}
    }
  \end{block}

  \vspace{0.3em}
  \begin{alertblock}{\footnotesize Salida de Consola — Parte 2: Impacto de Outliers}
    {\fontsize{6.5pt}{7.5pt}\selectfont
    \texttt{--- Impacto de Outliers en S vs IQR ---} \\
    \texttt{Sin outlier:  Desv. Estándar S = 3.45   | IQR = 4.75} \\
    \texttt{Con outlier:  Desv. Estándar S = 133.29 | IQR = 6.00} \\
    \texttt{Distorsión S: 38.6x veces mayor | Distorsión IQR: 1.3x (prácticamente estable)}
    }
  \end{alertblock}
\end{frame}

% Slide 10: Ejercicios del Curso - Nivel 1

% Slide 11: Ejercicios del Curso - Nivel 2

% Slide 12: Ejercicios del Curso - Nivel 3

% Slide 13: Ejercicios del Curso - Nivel 4

% Slide 14: Resumen de la Sección 1.3
\begin{frame}{Resumen de la Sección 1.3}
  \begin{block}{Lo que dominamos en esta sección}
    \begin{itemize}
      \item \textbf{Medida Clásica ($S^2$ y $S$):} Cuantifican la dispersión en torno al centro mediante desviaciones al cuadrado. La corrección de Bessel ($n-1$) garantiza un estimador insesgado. \pause
      \item \textbf{Normalización ($CV$):} El Coeficiente de Variación permite comparar la dispersión de variables con distintas escalas y unidades de medida. \pause
      \item \textbf{Robustez Intercuartílica ($IQR$):} El rango intercuartílico y la regla de Tukey ($1.5 \cdot IQR$) protegen nuestros modelos del impacto devastador de los atípicos extremos.
    \end{itemize}
  \end{block}
  \pause
  \vspace{0.6em}
  \begin{center}
    \color{TecAzulOscuro}\textbf{Pausa de Asimilación:} ¿Alguna duda sobre cuándo usar la corrección de Bessel ($n-1$) o cómo calcular el $IQR$?
    \color{black}
  \end{center}
\end{frame}

% Slide 15: Conexión con el Siguiente Capítulo
\begin{frame}{Conexión con el Próximo Capítulo}
  \begin{alertblock}{Siguiente Hito: Capítulo 02 — Probabilidad y Variables Aleatorias}
    Hasta ahora, hemos aprendido a resumir datos \emph{observados en el pasado} a través de tablas, medidas de tendencia central y de dispersión.
    \vspace{0.6em}
    \\ Para dar el salto hacia la \textbf{Estadística Inferencial} y el \textbf{Machine Learning predictivo}, necesitamos un lenguaje matemático para cuantificar el azar y la incertidumbre en eventos futuros:
    \begin{itemize}
      \item Axiomas y Fundamentos de Probabilidad
      \item Teorema de Bayes y Probabilidad Condicional
      \item Variables Aleatorias Discretas y Continuas
    \end{itemize}
  \end{alertblock}
\end{frame}

\end{document}
